% ============================================================
\chapter{Result and Discussion}
\label{ch:result-discussion}
% ============================================================

\section{Experimental Setup}
\label{sec:experimental-setup}

All runs were carried out on consumer hardware with training confined to the
CPU, using the pipeline described in
Section~\ref{sec:common-architecture}. Every run logs one row per generation to
a CSV file, and the figures in this chapter are drawn from those logs.

\noindent
Fitness values are not compared across approaches anywhere in this chapter. The
three approaches use different fitness functions with different scales and
different meanings, so a fitness of 500 under one is not a better or worse
result than a fitness of 500 under another. Where a comparison is made across
approaches it is made on distance covered, on network structure, or on the
consistency of a genome's performance across levels, all of which mean the same
thing regardless of how the genome was scored.

\noindent
Table~\ref{tab:run-summary} summarises the runs.

\begin{table}[H]
\centering
\caption{Summary of Training Runs}
\label{tab:run-summary}
\small
\begin{tabular}{@{}llrr@{}}
\toprule
\textbf{Approach} & \textbf{Run} & \textbf{Generations} & \textbf{Best distance} \\
\midrule
I   & Level 1-1        & 1102 & 3267 \\
I   & Level 2-1        &  744 & 2458 \\
I   & Level 3-1        & 4054 & 2962 \\
\addlinespace
II  & Multi-level      &  154 & 3267 \\
\addlinespace
III & Mean aggregation &  161 & --- \\
III & Spread penalty   &   87 & --- \\
\bottomrule
\end{tabular}
\end{table}

\noindent
Distance is not reported for Approach III because a moment run does not
measure it the same way: a genome starts partway through a level and is scored
on the ground it covers from there, so the figure is not comparable with a
full-level run.

% ============================================================
\section{Approach I: Level-Specific Training}
\label{sec:results-approach-1}

\subsection{Distance and Level Coverage}

Each level was trained with its own population. Figure~\ref{fig:a1-distance}
shows the furthest point reached over the course of each run, and
Figure~\ref{fig:a1-progress} expresses the same data as a percentage of the
level's length so the three runs can be read against one another.

\begin{figure}[H]
\centering
\includegraphics[width=0.85\textwidth]{results/a1_distance}
\caption{Best distance reached per generation, Approach I}
\label{fig:a1-distance}
\end{figure}

\begin{figure}[H]
\centering
\includegraphics[width=0.85\textwidth]{results/a1_progress}
\caption{Proportion of each level covered, Approach I}
\label{fig:a1-progress}
\end{figure}

\noindent
Level 1-1 was solved. The population reached the flagpole at generation 500 and
held it for the remaining 600 generations of the run. Level 3-1 reached 90.7\%
of the level, but took 4054 generations to do so and spent most of that time on
a plateau: it sat at roughly 74\% from generation 1500 to generation 3000
before a late improvement carried it most of the way. Level 2-1 was the
hardest, stalling at 75.3\% and not improving over the last 200 generations of
its run.

\noindent
The shape of all three curves is the same and is worth noting: progress is not
gradual. Long flat stretches are broken by sudden jumps, which is what
evolutionary search looks like when a structural mutation finally produces
behaviour that clears an obstacle the population has been stuck against. The
flat sections are not wasted time; they are the population accumulating
variation that has not yet paid off.

\subsection{Structural Complexity}

Figure~\ref{fig:a1-structure} shows how the champion network's structure
developed.

\begin{figure}[H]
\centering
\includegraphics[width=0.95\textwidth]{results/a1_structure}
\caption{Champion network structure over training, Approach I}
\label{fig:a1-structure}
\end{figure}

\noindent
This is NEAT complexifying as intended. Every run starts from the same place,
zero hidden nodes and 1248 enabled connections, since the initial genomes are
fully connected with no hidden layer. From there the two quantities move in
opposite directions. Hidden nodes accumulate, reaching 50 on Level 1-1 and 8 on
the other two runs. Enabled connections fall sharply, from 1248 down to around
100 on Level 1-1 and 250 to 320 on the others.

\noindent
The network is therefore not growing in the naive sense. It is discarding most
of the dense input-to-output wiring it started with and replacing a small part
of it with hidden structure. What survives is a much sparser network with a few
intermediate nodes, which is a more useful description of the policy than the
fully connected sheet it began as. Level 1-1, the run that actually solved its
level, is also the run that went furthest in both directions, adding the most
hidden nodes and pruning to the fewest connections.

\subsection{Learning Rate}

Figure~\ref{fig:cmp-improvement} normalises each run against its own starting
distance and plots it against the fraction of the run completed, which makes
the three learning curves comparable despite their very different lengths.

\begin{figure}[H]
\centering
\includegraphics[width=0.85\textwidth]{results/cmp_improvement}
\caption{Relative improvement over a training run, Approach I}
\label{fig:cmp-improvement}
\end{figure}

\noindent
All three runs gain most of their improvement early. By the time 20\% of each
run has elapsed, Level 1-1 has already improved 6.7-fold over its generation-0
distance and Level 2-1 5.8-fold, and the remaining 80\% of the run adds
comparatively little. This is the practical argument against simply training
for longer: Level 3-1 ran for 4054 generations, nearly four times Level 1-1's
run, and finished with a smaller relative gain.

\subsection{Limitation}

The limitation of this approach is the one the curves cannot show. Each of
these genomes was evolved against a single fixed layout and is scored only on
that layout, so nothing in training rewarded behaviour that transfers. The
Level 1-1 champion has learned the specific sequence of obstacles in Level 1-1,
not how to play Super Mario Bros., and it performs poorly when placed in any
other level. This is the same overfitting reported for
MarI/O~\cite{sethbling2015}, and it is what motivated the second approach.

% ============================================================
\section{Approach II: Multi-Level Training}
\label{sec:results-approach-2}

Approach II evaluates one genome across several levels and combines the results
into a single score, so that selection favours genomes that are competent
everywhere rather than genomes tuned to one layout. The run analysed here
resumes from generation 806 of an earlier run and continues to generation 959.

\begin{figure}[H]
\centering
\includegraphics[width=0.85\textwidth]{results/a2_distance}
\caption{Distance reached per generation, Approach II}
\label{fig:a2-distance}
\end{figure}

\noindent
Figure~\ref{fig:a2-distance} shows distance holding between roughly 2900 and
3267 across the window, with the running best reaching the full level length.
The per-generation trace varies considerably from one generation to the next
while the running best barely moves, which is the signature of a population
that has converged: the champion is stable and the variation is the rest of the
population exploring around it.

\begin{figure}[H]
\centering
\includegraphics[width=0.95\textwidth]{results/a2_structure}
\caption{Champion network structure, Approach II}
\label{fig:a2-structure}
\end{figure}

\noindent
The structural picture in Figure~\ref{fig:a2-structure} is markedly different
from Approach I. The champion carries around 150 hidden nodes, far more than
the 8 to 50 seen in the level-specific runs, while its enabled connections
continue to fall, from 304 to 235 over the window. A genome that has to handle
several levels ends up with substantially more internal structure than one
tuned to a single layout, which is what one would expect if the extra nodes are
encoding responses to situations that only arise on some of the levels.

\begin{figure}[H]
\centering
\includegraphics[width=0.85\textwidth]{results/a2_species}
\caption{Active species count, Approach II}
\label{fig:a2-species}
\end{figure}

\noindent
Figure~\ref{fig:a2-species} shows the species count moving between 1 and 12.
The periods where it falls to a very low value are the population converging
onto a single strategy, and the recoveries afterwards are the reset mechanism
described in Section~\ref{subsec:approach-multilevel} rebuilding diversity. The
mechanism is visibly doing its job here, though the fact that it has to
intervene repeatedly indicates how strongly this configuration tends toward
convergence.

% ============================================================
\section{Approach III: Moment-Based Training}
\label{sec:results-approach-3}

\subsection{Mean Aggregation}

The first moment-based scheme scores a genome as the plain mean of its
per-moment fitness across the whole batch. Figure~\ref{fig:a3-mean} plots the
champion's fitness on each of the four levels separately, which is what makes
the problem with this scheme visible.

\begin{figure}[H]
\centering
\includegraphics[width=0.85\textwidth]{results/a3_mean}
\caption{Per-level champion fitness under mean aggregation}
\label{fig:a3-mean}
\end{figure}

\noindent
The four curves separate and stay separated. Averaged over the last ten
generations of the run, the champion scores 190.2 on Level 1-1 and 189.2 on
Level 4-1, but only 88.5 on Level 2-1, a spread of 101.7 between its best and
worst level.

\noindent
This is exactly the failure the mean invites. A genome is selected on its
average, so a large gain on one level fully compensates a large loss on
another. Selection is therefore indifferent between a genome that plays all
four levels adequately and one that plays two of them well and a third badly,
provided the averages match. Since some levels are easier to improve on than
others, the population drifts toward the easy ones and lets the hard one
degrade, which is what the gap between the Level 1-1 and Level 2-1 curves
records.

\subsection{Spread Penalty}

The second scheme subtracts a multiple of the standard deviation across levels
from the mean, so a genome is rewarded for being good and for being evenly
good.

\begin{figure}[H]
\centering
\includegraphics[width=0.85\textwidth]{results/a3_spread}
\caption{Per-level champion fitness under the spread penalty}
\label{fig:a3-spread}
\end{figure}

\noindent
Figure~\ref{fig:a3-spread} shows the four curves running much closer together.
Over the last ten generations the champion scores 143.8, 125.7, 122.7 and 174.8
on Levels 1-1 through 4-1, a spread of 52.1 against the 101.7 seen under the
mean. The important part is where the change came from: Level 2-1, the weakest
level under the mean, improves from 88.5 to 125.7, while Level 1-1 falls from
190.2 to 143.8. The penalty has traded peak performance on the easy levels for
competence on the hard one, which is precisely what it was designed to do.

\subsection{Consistency}

Figure~\ref{fig:a3-consistency} states the comparison directly, plotting the
standard deviation of the champion's fitness across the four levels for both
schemes. Both curves are restricted to generations 51 to 137, the window the
two runs share, so they are compared over identical ground.

\begin{figure}[H]
\centering
\includegraphics[width=0.85\textwidth]{results/a3_consistency}
\caption{Consistency across levels under both aggregation schemes}
\label{fig:a3-consistency}
\end{figure}

\noindent
Over that window the mean scheme averages a cross-level standard deviation of
261, against 24 for the spread penalty, a reduction of roughly eleven times.
The mean scheme's curve is also erratic, climbing above 500 around generation
100 before dropping sharply, which reflects the champion changing to a genome
with a quite different balance across levels. The spread penalty's curve stays
flat and low for its entire run.

\noindent
The mean distance covered per moment supports the same conclusion. Under the
mean scheme it falls over the run, from 230 to 188, while under the spread
penalty it rises from 163 to 177. A scheme that was rewarding genuine
improvement should not see its distance measure decline, and the fact that the
mean scheme does is a further sign that it was optimising the score rather than
the behaviour.

\noindent
One qualification should be stated. The spread penalty run is shorter, 87
generations against 161, and it resumes from an existing checkpoint rather than
starting fresh, so the comparison is between the two schemes over a shared
window rather than between two complete runs from scratch. The size of the
difference, an order of magnitude in cross-level deviation, is large enough
that this does not change the conclusion, but a longer paired run would state
it more firmly.

% ============================================================
\section{Discussion}
\label{sec:discussion}

The three approaches form a sequence in which each one addresses the limitation
the previous one exposed.

\noindent
Approach I established that NEAT can evolve a competent agent for a fixed
level, solving Level 1-1 outright and covering most of Levels 2-1 and 3-1. It
also established the cost of doing so: those genomes are specialists, they do
not transfer, and one has to be trained per level. The result is a good
baseline and a poor agent.

\noindent
Approach II removed the per-level specialisation by requiring one genome to
play several levels before it is scored. This produced champions with
substantially more internal structure, consistent with a policy that has to
cover more situations, and the diversity trace shows the population repeatedly
converging and having to be rebuilt, which suggests the multi-level objective
is harder to satisfy than the single-level one.

\noindent
Approach III changed what a genome practises rather than what it is scored on,
evaluating from mid-level situations so that difficult sections are trained as
often as the opening stretch. Doing this exposed a problem that Approach II's
combined score had hidden: with a plain mean, consistency across levels is not
actually being selected for, and the population will quietly trade a hard level
away. The spread penalty corrects this and reduces cross-level deviation by
roughly eleven times, at a deliberate cost to peak performance on the easier
levels.

\noindent
The broader observation across all three is that the fitness function, not the
network, was where nearly every problem in this project was found and fixed.
Each of the failures documented here, the specialist that will not transfer,
the population that converges, and the mean that trades away a level, was a
case of the agent optimising exactly what it was asked to optimise. The
corrections were changes to the question, not to the learner.
